% Options for packages loaded elsewhere
% Options for packages loaded elsewhere
\PassOptionsToPackage{unicode}{hyperref}
\PassOptionsToPackage{hyphens}{url}
\PassOptionsToPackage{dvipsnames,svgnames,x11names}{xcolor}
%
\documentclass[
  english,
  review,
  number,
  sort,compress,
  authoryear,
  3p]{elsarticle}
\usepackage{xcolor}
\usepackage{amsmath,amssymb}
\setcounter{secnumdepth}{5}
\usepackage{iftex}
\ifPDFTeX
  \usepackage[T1]{fontenc}
  \usepackage[utf8]{inputenc}
  \usepackage{textcomp} % provide euro and other symbols
\else % if luatex or xetex
  \usepackage{unicode-math} % this also loads fontspec
  \defaultfontfeatures{Scale=MatchLowercase}
  \defaultfontfeatures[\rmfamily]{Ligatures=TeX,Scale=1}
\fi
\usepackage{lmodern}
\ifPDFTeX\else
  % xetex/luatex font selection
\fi
% Use upquote if available, for straight quotes in verbatim environments
\IfFileExists{upquote.sty}{\usepackage{upquote}}{}
\IfFileExists{microtype.sty}{% use microtype if available
  \usepackage[]{microtype}
  \UseMicrotypeSet[protrusion]{basicmath} % disable protrusion for tt fonts
}{}
\makeatletter
\@ifundefined{KOMAClassName}{% if non-KOMA class
  \IfFileExists{parskip.sty}{%
    \usepackage{parskip}
  }{% else
    \setlength{\parindent}{0pt}
    \setlength{\parskip}{6pt plus 2pt minus 1pt}}
}{% if KOMA class
  \KOMAoptions{parskip=half}}
\makeatother
% Make \paragraph and \subparagraph free-standing
\makeatletter
\ifx\paragraph\undefined\else
  \let\oldparagraph\paragraph
  \renewcommand{\paragraph}{
    \@ifstar
      \xxxParagraphStar
      \xxxParagraphNoStar
  }
  \newcommand{\xxxParagraphStar}[1]{\oldparagraph*{#1}\mbox{}}
  \newcommand{\xxxParagraphNoStar}[1]{\oldparagraph{#1}\mbox{}}
\fi
\ifx\subparagraph\undefined\else
  \let\oldsubparagraph\subparagraph
  \renewcommand{\subparagraph}{
    \@ifstar
      \xxxSubParagraphStar
      \xxxSubParagraphNoStar
  }
  \newcommand{\xxxSubParagraphStar}[1]{\oldsubparagraph*{#1}\mbox{}}
  \newcommand{\xxxSubParagraphNoStar}[1]{\oldsubparagraph{#1}\mbox{}}
\fi
\makeatother


\usepackage{longtable,booktabs,array}
\usepackage{calc} % for calculating minipage widths
% Correct order of tables after \paragraph or \subparagraph
\usepackage{etoolbox}
\makeatletter
\patchcmd\longtable{\par}{\if@noskipsec\mbox{}\fi\par}{}{}
\makeatother
% Allow footnotes in longtable head/foot
\IfFileExists{footnotehyper.sty}{\usepackage{footnotehyper}}{\usepackage{footnote}}
\makesavenoteenv{longtable}
\usepackage{graphicx}
\makeatletter
\newsavebox\pandoc@box
\newcommand*\pandocbounded[1]{% scales image to fit in text height/width
  \sbox\pandoc@box{#1}%
  \Gscale@div\@tempa{\textheight}{\dimexpr\ht\pandoc@box+\dp\pandoc@box\relax}%
  \Gscale@div\@tempb{\linewidth}{\wd\pandoc@box}%
  \ifdim\@tempb\p@<\@tempa\p@\let\@tempa\@tempb\fi% select the smaller of both
  \ifdim\@tempa\p@<\p@\scalebox{\@tempa}{\usebox\pandoc@box}%
  \else\usebox{\pandoc@box}%
  \fi%
}
% Set default figure placement to htbp
\def\fps@figure{htbp}
\makeatother



\ifLuaTeX
\usepackage[bidi=basic]{babel}
\else
\usepackage[bidi=default]{babel}
\fi
% get rid of language-specific shorthands (see #6817):
\let\LanguageShortHands\languageshorthands
\def\languageshorthands#1{}
\ifLuaTeX
  \usepackage[english]{selnolig} % disable illegal ligatures
\fi


\setlength{\emergencystretch}{3em} % prevent overfull lines

\providecommand{\tightlist}{%
  \setlength{\itemsep}{0pt}\setlength{\parskip}{0pt}}



 
\usepackage[]{natbib}
\bibliographystyle{elsarticle-harv}


\makeatletter
\@ifpackageloaded{caption}{}{\usepackage{caption}}
\AtBeginDocument{%
\ifdefined\contentsname
  \renewcommand*\contentsname{Table of contents}
\else
  \newcommand\contentsname{Table of contents}
\fi
\ifdefined\listfigurename
  \renewcommand*\listfigurename{List of Figures}
\else
  \newcommand\listfigurename{List of Figures}
\fi
\ifdefined\listtablename
  \renewcommand*\listtablename{List of Tables}
\else
  \newcommand\listtablename{List of Tables}
\fi
\ifdefined\figurename
  \renewcommand*\figurename{Figure}
\else
  \newcommand\figurename{Figure}
\fi
\ifdefined\tablename
  \renewcommand*\tablename{Table}
\else
  \newcommand\tablename{Table}
\fi
}
\@ifpackageloaded{float}{}{\usepackage{float}}
\floatstyle{ruled}
\@ifundefined{c@chapter}{\newfloat{codelisting}{h}{lop}}{\newfloat{codelisting}{h}{lop}[chapter]}
\floatname{codelisting}{Listing}
\newcommand*\listoflistings{\listof{codelisting}{List of Listings}}
\makeatother
\makeatletter
\makeatother
\makeatletter
\@ifpackageloaded{caption}{}{\usepackage{caption}}
\@ifpackageloaded{subcaption}{}{\usepackage{subcaption}}
\makeatother
\journal{Computers and Electronics in Agriculture}
\usepackage{bookmark}
\IfFileExists{xurl.sty}{\usepackage{xurl}}{} % add URL line breaks if available
\urlstyle{same}
\hypersetup{
  pdftitle={Supplementary Material for: Explainable machine learning for wheat above-ground biomass estimation under spatial cross-validation: sensor contributions and exploratory harvest forecasting},
  pdfauthor={Francisco Zambrano; Abel Herrera; Mauricio Molina-Rocco},
  pdflang={en},
  colorlinks=true,
  linkcolor={blue},
  filecolor={Maroon},
  citecolor={Blue},
  urlcolor={Blue},
  pdfcreator={LaTeX via pandoc}}


\setlength{\parindent}{6pt}
\begin{document}

\begin{frontmatter}
\title{Supplementary Material for: Explainable machine learning for
wheat above-ground biomass estimation under spatial cross-validation:
sensor contributions and exploratory harvest forecasting}
\author[1]{Francisco Zambrano%
%
}

\author[2]{Abel Herrera%
%
}

\author[3]{Mauricio Molina-Rocco%
%
}


\affiliation[1]{organization={Facultad de Economía, Negocios y Gobierno,
Universidad San Sebastián, Santiago, Chile},,postcodesep={}}
\affiliation[2]{organization={Hemera Centro de Observación de la Tierra,
Facultad de Ciencias, Ingeniería y Tecnología; Universidad Mayor,
Santiago, Chile},,postcodesep={}}
\affiliation[3]{organization={Laboratorio de Suelos y Sostenibilidad.
Departamento de Acuicultura y Recursos Agroalimentarios, Universidad de
Los Lagos, Osorno, 5311157, Chile},,postcodesep={}}

\cortext[cor1]{Corresponding author}



        





\end{frontmatter}
    

\section{Table S1: Additional vegetation
indices}\label{table-s1-additional-vegetation-indices}

Table 3 of the main text lists the 21 vegetation indices (VIs), derived
from Sentinel-2 (S2) and PlanetScope (PS) spectral bands, that were
retained as predictors after eliminating highly correlated indices
(Section 2.2.5 of the main text). Table S1 below lists the remaining 32
VIs that were part of the same 53-index base candidate set (Section
2.2.5) but were not retained in Table 3.

As described in the main text, a subset of these S2-derived indices were
additionally recalculated using the Sentinel-2 Band 8A (narrow
near-infrared) instead of Band 8 (broad near-infrared), following the
same formula with the near-infrared (NIR) term replaced by the Band 8A
reflectance; this is denoted with an asterisk (*) next to the S2 dataset
code, following the same convention as Table 3 of the main text.
Together, the 53 base indices (21 in Table 3 and 32 in Table S1) and
their 36 Band-8A-recalculated counterparts make up the full set of 89
candidate VIs referenced in Section 2.2.5 of the main text.

Band and index notation follows Table 3 of the main text: Red, Green,
Blue = visible bands; RE1, RE2, RE3 = S2 red-edge bands 5, 6, and 7; NIR
= near-infrared band (S2 Band 8 / PS Band 8); SWIR1, SWIR2 = S2
shortwave-infrared bands 11 and 12; Water vapour = S2 Band 9. Formula
definitions follow standard vegetation-index literature, including the
references cited in Section 2.2.5 of the main text for VI selection
\citep{Peng2023, Segarra2022, Uribeetxebarria2023, Wang2022, Wu2008, Zheng2017};
the exact implementation of every index listed here and in Table 3 is
available in the project code repository (script
\texttt{02\_calcular\_vis.R}).

\textbf{Table S1.} Additional vegetation indices (covariates) derived
from Sentinel-2 (S2) and PlanetScope (PS) spectral bands, computed as
part of the 53-index base candidate set (Section 2.2.5 of the main text)
but not retained in Table 3. Asterisk (*) next to the S2 dataset code
indicates that the index was also recalculated using Sentinel-2 Band 8A
(narrow NIR) instead of Band 8 (broad NIR).

\begin{longtable}[]{@{}
  >{\raggedright\arraybackslash}p{(\linewidth - 8\tabcolsep) * \real{0.0851}}
  >{\raggedright\arraybackslash}p{(\linewidth - 8\tabcolsep) * \real{0.2979}}
  >{\raggedright\arraybackslash}p{(\linewidth - 8\tabcolsep) * \real{0.1277}}
  >{\raggedright\arraybackslash}p{(\linewidth - 8\tabcolsep) * \real{0.2979}}
  >{\raggedright\arraybackslash}p{(\linewidth - 8\tabcolsep) * \real{0.1915}}@{}}
\toprule\noalign{}
\begin{minipage}[b]{\linewidth}\raggedright
N°
\end{minipage} & \begin{minipage}[b]{\linewidth}\raggedright
Dataset used
\end{minipage} & \begin{minipage}[b]{\linewidth}\raggedright
Name
\end{minipage} & \begin{minipage}[b]{\linewidth}\raggedright
Abbreviation
\end{minipage} & \begin{minipage}[b]{\linewidth}\raggedright
Formula
\end{minipage} \\
\midrule\noalign{}
\endhead
\bottomrule\noalign{}
\endlastfoot
1 & S2 & NDVI Variant (Red-edge 3) & NDVI\textsubscript{2} &
\((RE_{3} - Red)/(RE_{3} + Red)\) \\
2 & S2 & Red-edge NDVI (705 nm) & NDVI\textsubscript{705} &
\((RE_{2} - RE_{1})/(RE_{2} + RE_{1})\) \\
3 & S2*-PS & Modified Simple Ratio & MSR &
\((NIR/Red - 1)/\sqrt{(NIR/Red + 1)}\) \\
4 & S2 & Red-edge Modified Simple Ratio & MSR\textsubscript{705} &
\((RE_{2}/RE_{1} - 1)/\sqrt{(RE_{2}/RE_{1} + 1)}\) \\
5 & S2*-PS & Soil-Adjusted Vegetation Index & SAVI &
\(\frac{1.5 \cdot (NIR - Red)}{(NIR + Red + 0.5)}\) \\
6 & S2 & SAVI Variant (Red-edge 3) & SAVI\textsubscript{2} &
\(\frac{1.5 \cdot (RE_{3} - Red)}{(RE_{3} + Red + 0.5)}\) \\
7 & S2*-PS & Green SAVI & SAVI\textsubscript{green} &
\(\frac{1.5 \cdot (NIR - Green)}{(NIR + Green + 0.5)}\) \\
8 & S2*-PS & Green OSAVI & OSAVI\textsubscript{green} &
\(\frac{1.16 \cdot (NIR - Green)}{(NIR + Green + 0.16)}\) \\
9 & S2*-PS & Red-edge OSAVI & OSAVI\textsubscript{rededge} &
\(\frac{1.16 \cdot (NIR - RE_{1})}{(NIR + RE_{1} + 0.16)}\) \\
10 & S2 & - & TCARI/OSAVI\textsubscript{2} & \(TCARI/{OSAVI}_{2}\) \\
11 & S2* & SWIR 11 Related OSAVI & OSAVI\textsubscript{SWIR11} &
\(\frac{1.16 \cdot (NIR - SWIR_{1})}{(NIR + SWIR_{1} + 0.16)}\) \\
12 & S2 & SWIR 12 Related TCARI & TCARI\textsubscript{SWIR12} &
\(3 \cdot ((RE_{1} - SWIR_{2}) - 0.2 \cdot (RE_{1} - Green) \cdot (RE_{1}/SWIR_{2}))\) \\
13 & S2* & SWIR 12 Related OSAVI & OSAVI\textsubscript{SWIR12} &
\(\frac{1.16 \cdot (NIR - SWIR_{2})}{(NIR + SWIR_{2} + 0.16)}\) \\
14 & S2*-PS & Chlorophyll Index - Green & CI\textsubscript{green} &
\((NIR/Green) - 1\) \\
15 & S2*-PS & Chlorophyll Index - Red-edge & CI\textsubscript{rededge} &
\((NIR/RE_{1}) - 1\) \\
16 & S2*-PS & OSAVI x Chlorophyll Index Red-edge & OSAVI x
CI\textsubscript{rededge} &
\(\left[\frac{1.16 \cdot (NIR - Red)}{(NIR + Red + 0.16)}\right] \cdot \left[(NIR/RE_{1}) - 1\right]\) \\
17 & S2*-PS & Difference Vegetation Index & DVI & \(NIR - Red\) \\
18 & S2*-PS & Green DVI & DVI\textsubscript{green} & \(NIR - Green\) \\
19 & S2*-PS & Red-edge DVI & DVI\textsubscript{rededge} &
\(NIR - RE_{1}\) \\
20 & S2 & Two-Band Enhanced Vegetation Index & EVI2 &
\(\frac{2.5 \cdot (RE_{3} - Red)}{(1 + RE_{3} + 2.4 \cdot Red)}\) \\
21 & S2*-PS & Red-edge Vegetation Index 1 & REVI1 & \(NIR - RE_{1}\) \\
22 & S2* & Red-edge Vegetation Index 2 & REVI2 & \(NIR - RE_{2}\) \\
23 & S2*-PS & Weighted Difference Vegetation Index & WDVI &
\(NIR - 0.5 \cdot Red\) \\
24 & S2-PS & Green-Red Vegetation Index & GRVI &
\((Green - Red)/(Green + Red)\) \\
25 & S2*-PS & Green Ratio Vegetation Index & GRVI2 & \(NIR/Green\) \\
26 & S2*-PS & Green Atmospherically Resistant Index & GARI &
\(\frac{NIR - (Green - 1.7 \cdot (Blue - Red))}{NIR + (Green - 1.7 \cdot (Blue - Red))}\) \\
27 & S2*-PS & Leaf Chlorophyll Index & LCI &
\((NIR - RE_{1})/(NIR - Red)\) \\
28 & S2 & MERIS Terrestrial Chlorophyll Index & MTCI &
\((RE_{2} - RE_{1})/(RE_{1} - Red)\) \\
29 & S2* & Normalized Difference Red-edge Index 2 & NDRE2 &
\((NIR - RE_{2})/(NIR + RE_{2})\) \\
30 & S2 & Normalized Difference Red-edge Index (RE3/Green) & NDREI &
\((RE_{3} - Green)/(RE_{3} + Green)\) \\
31 & S2* & Normalized Difference Moisture Index & NDMI &
\((NIR - SWIR_{1})/(NIR + SWIR_{1})\) \\
32 & S2*-PS & Kernel NDVI & KNDVI &
\(\dfrac{1 - \exp\left(-\dfrac{(NIR-Red)^2}{2 \cdot 0.15^2}\right)}{1 + \exp\left(-\dfrac{(NIR-Red)^2}{2 \cdot 0.15^2}\right)}\) \\
\end{longtable}

\section{Equations S1-S3: Performance
metrics}\label{equations-s1-s3-performance-metrics}

Section 2.3.6 of the main text reports model performance using the mean
absolute error (MAE), root mean square error (RMSE), and coefficient of
determination (R²), defined in Equations S1-S3 below, where \(y_i\) and
\(\hat{y}_i\) are the observed and predicted AGB values for sample
\(i\), \(\bar{y}\) is the mean observed AGB across the \(n\) samples in
the evaluation set, and \(n\) is the number of samples.

\[
MAE = \frac{1}{n}\sum_{i=1}^{n}\left|y_i - \hat{y}_i\right| \tag{S1}
\]

\[
RMSE = \sqrt{\frac{1}{n}\sum_{i=1}^{n}\left(y_i - \hat{y}_i\right)^2} \tag{S2}
\]

\[
R^2 = 1 - \frac{\displaystyle\sum_{i=1}^{n}\left(y_i - \hat{y}_i\right)^2}{\displaystyle\sum_{i=1}^{n}\left(y_i - \bar{y}\right)^2} \tag{S3}
\]


\bibliography{references.bib}



\end{document}
